\section{주분석: 가격측도와 교차풀 제약}

\subsection{게시가격과 삼쌍승식 상태분포}

경주 $r$의 유효 출전마 집합을 $N_r$, 유효 출전마 수를 $n_r=|N_r|$라 하고,
순서 있는 상위 3마리의 상태공간을
\[
  \Omega_r=\{(i,j,k)\in N_r^3:i,j,k\text{는 서로 다르다}\}
\]
로 둔다. 삼쌍승식 상태 $(i,j,k)$의 게시 총지급배율을 $D^{\T}_{r,ijk}$라 하자.
KRA 배당률은 원금을 포함하므로 단위지급 청구권의 게시가격은
$x^{\T}_{r,ijk}=1/D^{\T}_{r,ijk}$다. 순수익 배당을 $O/1$로 표시하는 문헌의
가격 $1/(O+1)$과 같은 개념이지만, 표기법은 구분한다.

주분석에서는 역배당률을 경주와 풀 안에서 정규화한
\begin{equation}
  q^{\T}_{r,ijk}
  =\frac{x^{\T}_{r,ijk}}
  {\sum_{(a,b,c)\in\Omega_r}x^{\T}_{r,abc}}
  \label{eq:trifecta-state-price}
\end{equation}
을 삼쌍승식 정규화 가격측도라 부른다. 이상적인 패리뮤추얼 풀에서 이는 풀의
조합별 베팅금액 비중과 대응한다. 객관적 순위확률이나 Arrow--Debreu 상태가격을
복원한다고 전제하지 않는다. 정규화는 풀 전체에 공통으로 곱해지는 공제 요인을
제거하지만 조합별 반올림, 유동성 차이와 마감 직전 베팅의 영향은 제거하지 않는다.

\subsection{주변화 연산}

삼쌍승식 분포가 주어지면 단승식, 쌍승식, 복승식의 재구성 가격은 각각
\begin{align}
  \widehat p^{\W\leftarrow\T}_{r,i}
    &=\sum_{j\ne i}\sum_{k\notin\{i,j\}}q^{\T}_{r,ijk},
    \label{eq:win-marginal}\\
  \widehat p^{\Eact\leftarrow\T}_{r,ij}
    &=\sum_{k\notin\{i,j\}}q^{\T}_{r,ijk},
    \label{eq:exacta-marginal}\\
  \widehat p^{\Q\leftarrow\T}_{r,\{i,j\}}
    &=\sum_{k\notin\{i,j\}}
      \left(q^{\T}_{r,ijk}+q^{\T}_{r,jik}\right)
    \label{eq:quinella-marginal}
\end{align}
이다. 삼복승식은 세 말의 여섯 순열을 합한다.
\begin{equation}
  \widehat p^{\Trio\leftarrow\T}_{r,\{i,j,k\}}
  =\sum_{\pi\in S_3}q^{\T}_{r,\pi(i,j,k)}.
  \label{eq:trio-marginal}
\end{equation}

각 하위 승식 $m$의 게시가격도 같은 방식으로 정규화하여 $p^m_r$를 만든다.
$A^m_r$를 식 \eqref{eq:win-marginal}--\eqref{eq:trio-marginal}의 합산을 나타내는
0--1 incidence matrix라 하면 완전한 교차풀 제약은
\begin{equation}
  p^m_r=A^m_r q^{\T}_r
  \label{eq:cross-market-restriction}
\end{equation}
이다. 실제 분석은 등식의 성립을 전제하지 않고 좌우의 거리를 측정한다.

\subsection{공통 가격측도 해석의 유지가정}

식 \eqref{eq:cross-market-restriction}을 하나의 공통 가격측도에서 나온 주변분포로
해석하려면 네 가지 유지가정이 필요하다. 첫째, 각 풀의 최종 가격이 같은
정보시점을 반영하고 취소마와 지급사건의 정의가 일치해야 한다. 둘째, 공제율은
풀 안에서 공통 배율로 작용하고 반올림·절사·적중자 없음 규칙이 조합별
상대가격을 크게 바꾸지 않아야 한다. 셋째, 별도 풀의 참가자 구성과 대표적
신념·평가기능이 같거나 안정적으로 연결되어야 한다. 넷째, 유동성과 최소
베팅단위가 풀별 비선형 왜곡을 만들지 않아야 한다.

이 가정이 모두 성립하지 않아도 주변화 연산 자체는 정확하다. 다만 이때 거리는
서로 다른 풀 가격측도의 \emph{정적 정렬 정도}이며, 무차익 상태가격의 존재나
개별 참가자의 합리성을 검정하지 않는다. 이 구분 때문에 본 논문은
``공통 상태가격을 입증한다''보다 ``공통 가격측도와의 양립 정도를 측정한다''고
표현한다.

\subsection{Harville 기준모형}

단승식 정규화 가격 $p^{\W}_{r,i}$만을 이용하는 Harville형 순차
모형\citep{harville1973}은
\begin{equation}
  q^{H}_{r,ijk}
  =p^{\W}_{r,i}
   \frac{p^{\W}_{r,j}}{1-p^{\W}_{r,i}}
   \frac{p^{\W}_{r,k}}{1-p^{\W}_{r,i}-p^{\W}_{r,j}}
  \label{eq:harville}
\end{equation}
로 정의한다. 이는 한 말이 먼저 들어온 뒤 남은 말의 상대확률이 비례적으로
재조정된다는 강한 가정을 둔다. 삼쌍승식 가격측도의 주변분포가
$A^m_rq^H_r$보다 실제 쌍승·복승·삼복승 가격에 더 가깝다면, 그 차이를
삼쌍승식 가격이 담은 순위 의존 구조의 증분으로 해석한다. 두 모형 모두
시장가격에서 출발하므로 이는 인과적 정보효과가 아니라 가격 설명력의 차이다.
